\documentclass[a4paper,12pt]{article}
\usepackage[T2A]{fontenc}
\usepackage[utf8]{inputenc}
\usepackage[russian]{babel}
\usepackage{amsmath,amssymb,icomma}
\usepackage[a4paper,left=20mm,right=20mm,top=20mm,bottom=22mm,headheight=25pt,headsep=8mm]{geometry}
\usepackage{microtype}
\usepackage{xcolor,tikz}
\usetikzlibrary{arrows.meta,positioning,shapes.geometric,calc}
\usepackage{booktabs,tabularx,enumitem,fancyhdr,titlesec}
\definecolor{accent}{HTML}{008080}
\definecolor{darkaccent}{HTML}{004D4D}
\definecolor{textgray}{HTML}{333333}
\color{textgray}
\setlength{\parindent}{0pt}
\setlength{\parskip}{6pt}
\setlength{\emergencystretch}{2em}
\raggedbottom
\titleformat{\section}{\large\color{darkaccent}}{\thesection.}{0.5em}{}[\vspace{2pt}\color{accent}\titlerule]
\titlespacing*{\section}{0pt}{14pt}{9pt}
\newcommand{\example}[1]{\par\medskip\textcolor{darkaccent}{\textbf{#1}}\par}
\newcommand{\checkitem}{\item[$\square$]}
\newcommand{\prob}{\mathop{\mathrm{P}}}
\pagestyle{fancy}
\fancyhf{}
\fancyhead[L]{\small Теория вероятностей}
\fancyhead[R]{\small Марк Гукин}
\fancyfoot[C]{\small\thepage}
\renewcommand{\headrulewidth}{0.35pt}
\renewcommand{\headrule}{\hbox to\headwidth{\color{accent}\leaders\hrule height \headrulewidth\hfill}}
\setlist[enumerate]{leftmargin=*,itemsep=8pt,topsep=4pt}
\setlist[itemize]{leftmargin=*,itemsep=3pt,topsep=3pt}
\begin{document}
\begin{titlepage}
\centering
\vspace*{30mm}
{\Large\color{darkaccent} МАТЕМАТИКА\par}
\vspace{16mm}
{\huge Чек-лист\par}
\vspace{10mm}
{\LARGE\color{accent} Теория вероятностей\par}
\vspace{4mm}
{\large Случайные события и многоэтапные испытания\par}
\vspace{16mm}
{\large Подготовка к ЕГЭ по профильной математике\par}
\vfill
{\large Марк Гукин\par}
\vspace{11mm}
25 ноября 2026 г.\par
\vspace{18mm}
{\small Лёвин Артём Александрович\par
эксклюзивно для Марка Гукина\par}
\vspace{18mm}
\begin{tikzpicture}
\draw[accent,line width=0.8pt] (0,0) .. controls (2,0.8) and (4,-0.8) .. (6,0) .. controls (8,0.8) and (10,-0.8) .. (12,0);
\fill[accent] (6,0) circle (1.2pt);
\end{tikzpicture}
\vspace{6mm}
\end{titlepage}
\setcounter{page}{1}
\section{Равновозможные исходы}
Случайный опыт — действие, результат которого заранее неизвестен. Элементарный исход — один конкретный результат опыта. Событие $A$ объединяет все исходы, при которых выполняется нужное условие.

Если все $n$ элементарных исходов равновозможны и $m$ из них благоприятствуют событию $A$, то
\[ \prob(A)=\frac{m}{n},\qquad 0\leq \prob(A)\leq 1. \]
Невозможное событие имеет вероятность $0$, достоверное — $1$.

\example{Пример 1. Случайный выбор}
В соревновании участвуют 10 спортсменов из России, 4 из Беларуси, 8 из Казахстана и 3 из Украины. Порядок выступлений определяется жеребьёвкой. Вероятность, что последним выступит спортсмен из Казахстана:
\[ n=10+4+8+3=25,\qquad \prob(A)=\frac{8}{25}=0,32. \]
При подсчёте учитывайте все объекты, среди которых проводится выбор.

\example{Пример 2. Нечётная сумма на двух костях}
У каждой честной кости по 3 чётных и 3 нечётных значения. При независимых бросках всего $6\cdot6=36$ упорядоченных исходов. Нечётная сумма получается в двух случаях: чётное + нечётное или нечётное + чётное.
\[ m=3\cdot3+3\cdot3=18,\qquad \prob(A)=\frac{18}{36}=0,5. \]
Для трёх костей нечётная сумма означает одну или три нечётные грани. Получаем четыре схемы чётности, по $3^3$ исходов в каждой:
\[ \prob(A)=\frac{4\cdot3^3}{6^3}=0,5. \]

\example{Пример 3. Фиксированное испытание}
Честную монету бросают независимо трижды. Требуется орёл при втором броске. Первый и третий броски могут иметь любой результат:
\[ \prob(A)=\frac{2\cdot1\cdot2}{2^3}=\frac12. \]

\example{Геометрическая вероятность: часы}
Если момент остановки часов равномерно распределён по 12-часовому циклу, промежуток от 11 до 1 по ходу стрелки занимает 2 часа. Вероятность остановки стрелки на этой дуге равна $2/12=1/6$.

\clearpage
\section{Сложение и умножение вероятностей}
Пусть $A$ и $B$ — события. Запись $A\cap B$ означает, что произошли оба события; $A\cup B$ — что произошло хотя бы одно из них.

\begin{tabularx}{\linewidth}{@{}p{0.32\linewidth}X@{}}
\toprule
Условие & Правило \\\midrule
Независимые события & $\prob(A\cap B)=\prob(A)\prob(B)$ \\
Несовместные события & $\prob(A\cup B)=\prob(A)+\prob(B)$ \\
Произвольные события & $\prob(A\cup B)=\prob(A)+\prob(B)-\prob(A\cap B)$ \\
Противоположное событие & $\prob(\overline A)=1-\prob(A)$ \\
\bottomrule
\end{tabularx}

Независимость означает, что наступление одного события не меняет вероятность другого (при условии, что первое имеет ненулевую вероятность). Несовместные события не могут произойти одновременно.

\example{Пример 4. Ни один автобус не свободен}
Вероятность, что каждый из двух автобусов свободен, равна $0,55$. При независимости их состояний вероятность того, что оба заняты:
\[ (1-0,55)^2=0,45^2=0,2025. \]

\example{Пример 5. Заданный порядок результатов}
Вероятность попадания при каждом независимом выстреле равна $0,7$. Первые четыре выстрела должны быть удачными, последний — неудачным.
\[ \prob(A)=0,7^4\cdot0,3=0,07203\approx0,07. \]
Порядок задан, поэтому учитывается ровно одна последовательность.

\example{Пример 6. Разность вложенных событий}
Вероятность, что прибор прослужит больше года, равна $0,96$, а больше двух лет — $0,91$. Событие «больше двух лет» входит в событие «больше года». Поэтому
\[ \prob(1<T\leq 2)=0,96-0,91=0,05. \]
Для непрерывного срока службы вероятность ровно двух лет равна нулю, поэтому результат совпадает с вероятностью $1<T<2$.

\clearpage
\section{Хотя бы один: противоположное событие}
Если событие включает много подходящих комбинаций, сначала проверьте, удобно ли посчитать вероятность его дополнения.
\[ \boxed{\prob(\text{хотя бы одно})=1-\prob(\text{ни одного})} \]
Для независимых испытаний с вероятностями успеха $p_1,\ldots,p_n$:
\[ \prob(\text{хотя бы один успех})=1-(1-p_1)\cdots(1-p_n). \]
Если все вероятности одинаковы, формула принимает вид $1-(1-p)^n$.

\example{Пример 7. Исправность батареек}
Вероятность брака одной батарейки равна $0,06$. Две батарейки выбираются так, что их состояния независимы. Обозначим исправность знаком $+$, брак — знаком $-$.
\begin{center}
\renewcommand{\arraystretch}{1.4}
\begin{tabular}{@{}c c c@{}}
\toprule
Первая / вторая & $+$ ($0,94$) & $-$ ($0,06$) \\\midrule
$+$ ($0,94$) & $0,8836$ & $0,0564$ \\
$-$ ($0,06$) & $0,0564$ & $0,0036$ \\
\bottomrule
\end{tabular}
\end{center}
Все четыре клетки несовместны; сумма их вероятностей равна единице.
\begin{align*}
\prob(\text{хотя бы одна неисправна})&=1-0,94^2=0,1164,\\
\prob(\text{хотя бы одна исправна})&=1-0,06^2=0,9964,\\
\prob(\text{ровно одна исправна})&=2\cdot0,94\cdot0,06=0,1128.
\end{align*}

\example{Пример 8. Три платёжных автомата}
Каждый из трёх автоматов неисправен с вероятностью $0,1$. Если состояния независимы, вероятность того, что хотя бы один автомат исправен:
\[ 1-0,1^3=0,999. \]

\textcolor{darkaccent}{\textbf{Проверяйте перевод условия:}}
\begin{itemize}
\checkitem «Хотя бы один исправен» — дополнение к «все неисправны».
\checkitem «Хотя бы один неисправен» — дополнение к «все исправны».
\checkitem «Ровно один» — только один успех; необходимо учитывать все возможные положения этого успеха.
\end{itemize}

\clearpage
\section{Дерево вероятностей и полная вероятность}
Дерево удобно, когда случайный опыт состоит из нескольких этапов. Каждая ветвь показывает возможный результат этапа и его вероятность. На каждом разветвлении сумма условных вероятностей исходящих ветвей равна $1$.

\textcolor{darkaccent}{\textbf{Правило:}} вдоль пути перемножайте вероятности ветвей; вероятности различных подходящих конечных путей складывайте. Вероятность на очередной ветви учитывает результаты предыдущих этапов.

\example{Пример 9. Контроль качества батареек}
Доля неисправных батареек равна $0,02$. Контроль бракует неисправную батарейку с вероятностью $0,99$, исправную — с вероятностью $0,01$. Найдём вероятность браковки.
\begin{center}
\begin{tikzpicture}[x=1cm,y=1cm,>=Stealth, every node/.style={font=\small,align=center},edge/.style={draw=accent,thick,-{Stealth[length=2mm]}}]
\node (s) at (0,0) {Батарейка};
\node (g) at (5,1.8) {Исправна};
\node (b) at (5,-1.8) {Неисправна};
\node[text=darkaccent] (gr) at (11,2.8) {Бракована};
\node (ga) at (11,0.8) {Принята};
\node[text=darkaccent] (br) at (11,-0.8) {Бракована};
\node (ba) at (11,-2.8) {Принята};
\draw[edge] (s) -- node[above,sloped] {$0,98$} (g);
\draw[edge] (s) -- node[below,sloped] {$0,02$} (b);
\draw[edge] (g) -- node[above,sloped] {$0,01$} (gr);
\draw[edge] (g) -- node[below,sloped] {$0,99$} (ga);
\draw[edge] (b) -- node[above,sloped] {$0,99$} (br);
\draw[edge] (b) -- node[below,sloped] {$0,01$} (ba);
\end{tikzpicture}
\end{center}
Подходят два несовместных пути, которые оканчиваются браковкой:
\[ \prob(A)=0,98\cdot0,01+0,02\cdot0,99=0,0296. \]
В общем виде, если $H$ и $\overline H$ — два противоположных случая, то
\[ \prob(A)=\prob(H)\prob(A\mid H)+\prob(\overline H)\prob(A\mid\overline H). \]
Запись $\prob(A\mid H)$ обозначает вероятность $A$ при условии, что $H$ произошло. Формула выше применяется, когда обе условные вероятности определены.

\example{Пример 10. Выбор пути}
На трёх последовательных развилках нужная ветвь выбирается с вероятностью $1/2$ при условии достижения соответствующей развилки. Вероятность пройти весь заданный маршрут:
\[ \frac12\cdot\frac12\cdot\frac12=\frac18. \]

\clearpage
\section{Таблицы исходов и испытания с остановкой}
\example{Пример 11. Две футбольные игры}
За победу дают 3 очка, за ничью — 1, за поражение — 0. Вероятность ничьей $0,1$, вероятности победы и поражения одинаковы, поэтому каждая равна $(1-0,1)/2=0,45$. Для выхода в следующий круг нужно набрать не менее 4 очков в двух играх. Считаем результаты игр независимыми.
\begin{center}
\renewcommand{\arraystretch}{1.4}
\begin{tabular}{@{}c c c c@{}}
\toprule
Первая / вторая & 3 ($0,45$) & 1 ($0,1$) & 0 ($0,45$) \\\midrule
3 ($0,45$) & \textcolor{darkaccent}{\textbf{6}} & \textcolor{darkaccent}{\textbf{4}} & 3 \\
1 ($0,1$) & \textcolor{darkaccent}{\textbf{4}} & 2 & 1 \\
0 ($0,45$) & 3 & 1 & 0 \\
\bottomrule
\end{tabular}
\end{center}
Подходят три выделенные клетки: $(3,3)$, $(3,1)$, $(1,3)$. Вероятность не выйти в следующий круг:
\[ 1-(0,45^2+2\cdot0,45\cdot0,1)=0,7075\approx0,71. \]

\example{Пример 12. Повторный анализ только при отрицательном первом}
Вероятность ложного положительного результата у здоровой собаки равна $0,005$. Если первый анализ отрицательный, делают второй; третий не назначается. Предположим, что вероятность ложного положительного результата второго анализа при отрицательном первом также равна $0,005$.

Ложный диагноз возможен по двум несовместным путям: первый анализ положителен, либо первый отрицателен и второй положителен.
\[ \prob(A)=0,005+0,995\cdot0,005=0,009975\approx0,01. \]
После положительного первого результата испытание прекращается. Не добавляйте множитель вероятности второго анализа к этому пути.

\textcolor{darkaccent}{\textbf{Контроль перед ответом}}
\begin{itemize}
\checkitem Проверьте, что перечислены все подходящие исходы без повторений.
\checkitem Уточните, равновозможны ли исходы и независимы ли испытания.
\checkitem Проверьте, что результат находится в диапазоне от $0$ до $1$.
\checkitem Округляйте только окончательный результат по условию задачи.
\end{itemize}

\clearpage
\section*{Алгоритм: выбор способа решения}
\vfill
\begin{center}
\begin{tikzpicture}[font=\small,>=Stealth,
 terminal/.style={ellipse,draw=accent,thick,align=center,minimum width=4.2cm,minimum height=0.85cm},
 action/.style={rectangle,rounded corners=2pt,draw=accent,align=center,text width=5.4cm,minimum height=1cm,inner sep=6pt},
 decision/.style={diamond,draw=accent,aspect=2.4,align=center,text width=3.5cm,inner sep=2pt},
 line/.style={-{Stealth[length=2mm]},draw=darkaccent,thick}]
\node[terminal] (start) at (0,0) {Начало};
\node[action] (event) at (0,-1.7) {Запишите искомое событие.\\ Выделите этапы опыта.};
\node[decision] (comp) at (0,-4.0) {Дополнение\\ проще?};
\node[action,text width=3.8cm] (other) at (5.5,-4.0) {Замените событие\\ противоположным.};
\node[decision] (equal) at (0,-7.2) {Равновозможные\\ исходы или\\ равномерная дуга?};
\node[action,text width=4.0cm] (count) at (-4.4,-10.4) {Дискретный опыт: $m/n$.\\ Равномерная дуга: отношение длин.};
\node[action,text width=5.5cm] (tree) at (3.7,-10.4) {Составьте таблицу или дерево.\\ Умножьте вероятности вдоль пути; сложите подходящие несовместные пути.};
\node[action] (back) at (0,-13.6) {Если считали дополнение,\\ вычтите результат из $1$.};
\node[action] (verify) at (0,-15.5) {Проверьте $0\leq P\leq1$;\\ округлите по условию.};
\node[terminal] (finish) at (0,-17.3) {Ответ};
\draw[line] (start)--(event);
\draw[line] (event)--(comp);
\draw[line] (comp.east)--node[above]{да}(other.west);
\coordinate (join) at ($(equal.north)+(0,0.45)$);
\draw[draw=darkaccent,thick] (comp.south)--node[left]{нет}(join);
\draw[draw=darkaccent,thick] (other.south) |- (join);
\fill[darkaccent] (join) circle (1.2pt);
\draw[line] (join)--(equal.north);
\draw[line] (equal.west) -| node[pos=0.25,above]{да}(count.north);
\draw[line] (equal.east) -| node[pos=0.25,above]{нет}(tree.north);
\draw[line] (count.south) |- (back.west);
\draw[line] (tree.south) |- (back.east);
\draw[line] (back)--(verify);
\draw[line] (verify)--(finish);
\end{tikzpicture}
\end{center}
\vfill
\clearpage
\section*{Тренировочный блок}
Выполните 10 заданий самостоятельно. Записывайте событие, способ подсчёта и итоговую вероятность. Промежуточные вычисления не округляйте.

\textcolor{darkaccent}{\textbf{Базовый уровень}}
\begin{enumerate}
\item В партии 800 флеш-накопителей 12 неисправны. Покупатель случайным образом выбирает один накопитель. Найдите вероятность того, что он исправен.
\item Три честные игральные кости бросают независимо. Найдите вероятность того, что сумма выпавших очков окажется чётной.
\item Часы остановились в случайный момент, равномерно распределённый по 12-часовому циклу. Найдите вероятность того, что часовая стрелка достигла отметки 10, но ещё не дошла до отметки 1 по ходу движения.
\item Вероятность того, что принтер проработает больше одного года, равна $0,92$, больше двух лет — $0,84$. Найдите вероятность того, что он проработает больше года, но не больше двух лет.
\end{enumerate}

\textcolor{darkaccent}{\textbf{Средний уровень}}
\begin{enumerate}[start=5]
\item Стрелок делает четыре независимых выстрела. Вероятность попадания при каждом выстреле равна $0,8$. Найдите вероятность того, что первые три выстрела будут удачными, а последний — неудачным.
\item В диспетчерской два независимых канала связи. Первый свободен с вероятностью $0,65$, второй — с вероятностью $0,6$. Найдите вероятность того, что ни один канал не свободен.
\item У каждого из трёх датчиков вероятность неисправности равна $0,04$. Состояния датчиков независимы. Найдите вероятность того, что хотя бы один датчик неисправен.
\end{enumerate}

\clearpage
\section*{Тренировочный блок (продолжение)}
\textcolor{darkaccent}{\textbf{Повышенный уровень}}
\begin{enumerate}[start=8]
\item Команда проводит две независимые игры. За победу начисляется 3 очка, за ничью — 1, за поражение — 0. Вероятность ничьей в каждой игре равна $0,2$, вероятности победы и поражения одинаковы. Найдите вероятность того, что команда наберёт меньше 4 очков за две игры.
\vspace{17mm}
\item Автоматическая линия выпускает детали. Вероятность того, что деталь дефектна, равна $0,03$. Контроль бракует дефектную деталь с вероятностью $0,98$ и по ошибке бракует годную деталь с вероятностью $0,02$. Найдите вероятность того, что случайно выбранная деталь будет забракована.
\vspace{17mm}
\item При проверке исправного устройства вероятность ложного сигнала о неисправности равна $0,015$. Если первый тест не дал сигнала, назначается второй; в противном случае проверка прекращается. Вероятность ложного сигнала во втором тесте при отрицательном первом также равна $0,015$. Найдите вероятность получить хотя бы один ложный сигнал за всю процедуру. Ответ округлите до сотых.
\end{enumerate}

\clearpage
\section*{Ответы}
\begin{center}
\renewcommand{\arraystretch}{1.5}
\begin{tabularx}{0.94\linewidth}{@{}c X l@{}}
\toprule
№ & Проверяемый навык & Ответ \\\midrule
1 & Случайный выбор & $0,985$ \\
2 & Чётность суммы & $0,5$ \\
3 & Геометрическая вероятность & $0,25$ \\
4 & Вложенные события & $0,08$ \\
5 & Заданный порядок испытаний & $0,1024$ \\
6 & Совместное наступление независимых событий & $0,14$ \\
7 & Противоположное событие & $0,115264$ \\
8 & Таблица исходов & $0,68$ \\
9 & Полная вероятность & $0,0488$ \\
10 & Испытания с остановкой & $0,03$ \\
\bottomrule
\end{tabularx}
\end{center}
\vspace{8mm}
\textcolor{darkaccent}{\large Итоговая памятка}
\begin{center}
\renewcommand{\arraystretch}{1.35}
\begin{tabularx}{\linewidth}{@{}p{0.39\linewidth}X@{}}
\toprule
Сигнал в задаче & Полезное действие \\\midrule
Равновозможные исходы & Разделите число благоприятных исходов на общее число исходов. \\
Несколько независимых событий одновременно & Перемножьте их вероятности. \\
Несколько несовместных случаев & Сложите их вероятности. \\
Хотя бы один & Проверьте, проще ли посчитать противоположное событие. \\
Два вложенных события & Вычтите вероятность меньшего события из вероятности большего. \\
Несколько этапов, разные условия & Постройте дерево и подпишите условные вероятности. \\
Пары результатов, сумма очков & Составьте таблицу упорядоченных исходов. \\
\bottomrule
\end{tabularx}
\end{center}
\end{document}
